% typeset: Pdftex
% Afterwards compile with pdflatex > bibtex > pdflatex > pdflatex.
% in TeXShop preferences, changed edit from bibtex to biber
% beamer likes biber
% latex likes bibtex

% https://tex.stackexchange.com/questions/270633/beamer-and-the-pause-command
% https://tex.stackexchange.com/questions/1423/is-there-a-nice-way-to-compile-a-beamer-presentation-without-the-pauses
%\documentclass[handout]{beamer}
%
\documentclass[]{beamer}

% point to content libraries
\newcommand{\pGlobal}			{../../../global/}
\newcommand{\pGlobalSetup}		{\pGlobal setup-global/}	
% \usepackage[]{xwatermark} option clash with hyperref
% \newwatermark[allpages,color=red!50,angle=45,scale=3,xpos=0,ypos=0]{DRAFT}
%\usepackage[printwatermark]		{xwatermark}
%	\newwatermark*[ firstpage, color = red!80, angle = 45, scale = 2.5, xpos = 0, ypos = 0 ]{\transparent{0.2}{DRAFT}}

\setbeamercovered{ transparent = 10 }
\usepackage{background}
%\newwatermark[firstpage,color=red!30, angle=45, scale=3, xpos=0, ypos=0]{\transparent{0.2}{DRAFT}}
\usepackage{transparent}
%\usepackage{multimedia} https://tex.stackexchange.com/questions/485385/playing-mp4-file-inside-beamer

% load LaTeX packages and personal macros
\input{\pGlobalSetup setup-global}
\input{\pLocalSetup setup-local}
%\usepackage[backend=biber, bibstyle=numeric, citestyle=numeric, sorting=none]{biblatex}
\usepackage[ backend = biber ]{biblatex}
\bibliography{\pSections seismic.bib}

\DeclareMathOperator*{\argmin}{arg\,min}
\newcommand{\normt}[1]			{\left\lVert #1 \right\rVert_{2}}
\newcommand{\normts}[1]			{\left\lVert #1 \right\rVert_{2}^{2}}

\title
[Seismic Data: Measurements and Analysis]
{Seismic Data\\Measurements and Analysis}

\author[DTRIAC Tutorial]{\Topa \\ \TopaEmail}
\institute{\dtriac} 
\medskip

\date{\today}
%\date{2018-08-31}

\begin{document}

\begin{frame}
	\titlepage
	\distA
\end{frame}

%\input{\pSections "sec splash"}

\begin{frame}\frametitle{Underground Explosion A}
\center
	\href{https://apps.dtic.mil/sti/pdfs/ADA627744.pdf\#page=40}{
	\begin{overpic}[ scale = 0.30 ]
		{\pLocalGraphics ug-shot-01}
		%\put(50, 50)	{\colorbox{white}{$a+b$}}
	\end{overpic}}
\end{frame}

\begin{frame}\frametitle{Overview}
	\tableofcontents[hideallsubsections]
\end{frame}

%
% main sections
\input{\pSections "sec-seismology"}
\input{\pSections "sec-gtune"}
\input{\pSections "sec-inverse"}

{\tiny{
%\begin{frame}[allowframebreaks,shrink=10]\frametitle{Bibliography}
\begin{frame}[allowframebreaks]\frametitle{Bibliography}
	\nocite{*}
	\printbibliography
\end{frame}}}
%
%\begin{frame}\frametitle{Physics of Shock Waves and High-Temperature Hydrodynamic Phenomena}
%\begin{enumerate}
%	\item \href{https://books.google.co.ck/books?id=PULCAgAAQBAJ&printsec=frontcover\#v=onepage&q&f=false}{Google Books}
%	\item \href{https://apps.dtic.mil/sti/citations/AD0622356}{DTIC}
%	\item \href{https://store.doverpublications.com/0486420027.html}{Dover Publishing}
%	\item \href{https://archive.org/details/physicsofshoakwa0000zeld}{Internet Archive}
%\end{enumerate}
%\end{frame}
% physics-of-shock-waves-and-high-temperature-hydrodynamic-phenomena-18.jpeg

%

% info slide during questions
\begin{frame}
	\titlepage
\end{frame}
%

\end{document}

%\tiny
%\scriptsize
%\footnotesize
%\small
%\normalsize
%\large
%\Large
%\LARGE
%\huge
%\Huge

%\, thin space (normally 1/6 of a quad);
%\> medium space (normally 2/9 of a quad);
%\; thick space (normally 5/18 of a quad);

\begin{frame}\frametitle{Frame Title}
\begin{enumerate}
	\item 
	\item 
	\item 
\end{enumerate}
\end{frame}

\begin{frame}\frametitle{Frame Title}
\begin{equation}
	\begin{array}{ccc} 
			%
			%
			%
	\end{array}
%\label{eq:}
\end{equation}
\end{frame}

\begin{frame}\frametitle{ }
\center
	\href{url}{
	\begin{overpic}[ scale = 1.0 ]
	{\pLocalGraphics graphic-file}
		%\put(-7,-10){Auxiliary text.}
	\end{overpic}}
\end{frame}

\begin{frame}\frametitle{Frame Title}
\begin{table}[htp]
%\caption{default}
\begin{center}
	\begin{tabular}{cc}
		%
		%
		%
	\end{tabular}
\end{center}
%\label{tab:label}
\end{table}%
\end{frame}

